\chapter{Practical Part}

\section{Overview of the Proposed System}

This section should provide a practical overview of the implemented diploma prototype. Describe the system as an end-to-end solution that includes an IoT data source, backend API, PostgreSQL database, dashboard frontend, rule-based condition scoring, and a supervised machine learning classifier.

You can briefly summarize the implemented technology stack and state that the prototype supports the full flow from telemetry ingestion to visualization, recommendations, and alerts.

\section{System Architecture}

This section should describe the architecture of the implemented system in layers and explain the role of each subsystem.

\subsection{Client-side}

In this subsection, describe the web interface implemented for users and administrators. Explain that the client side provides authentication pages, plant views, sensor management, analytics charts, alert handling, and admin monitoring.

\subsection{Server-side}

This subsection should describe the backend API layer, application services, repositories, schemas, authentication, and alert logic. Explain how the server processes user requests and telemetry, persists records, and exposes structured responses to the frontend.

\subsection{IoT and Sensor Data Flow}

This subsection should explain how the simulated or physical sensor sends telemetry to the backend ingest endpoint, how device authentication is performed, and how data moves into storage and user-facing monitoring views.

\begin{figure}[H]
    \centering
    \fbox{\rule{0pt}{2.2in}\rule{0.85\textwidth}{0pt}}
    \caption{Placeholder architecture diagram for the complete IoT-to-dashboard data flow}
    \label{fig:architecture-placeholder}
\end{figure}

Figure~\ref{fig:architecture-placeholder} is a sample figure block. Replace the framed placeholder with the final architecture image exported from a diagram tool or prepared in vector format.

\subsection{AI/ML Condition Analysis Flow}

This subsection should describe the hybrid analysis flow implemented in the final system. Explain that the rule-based layer remains primary, while the Random Forest classifier provides supporting condition classification with confidence and class probabilities.

Suggested sequence to describe:
\begin{enumerate}
    \item sensor data is received and stored;
    \item current values and recent history are selected;
    \item rule-based condition scoring is calculated;
    \item trend features are derived;
    \item machine learning prediction is executed if the model is available;
    \item the dashboard returns a hybrid result with explanation and AI support fields.
\end{enumerate}

\section{Database Design}

This section should explain the structure of the PostgreSQL database and how the main entities support the implemented workflows.

\subsection{Entity-Relationship Model}

In this subsection, describe the key entities and their relationships, for example users, plants, sensors, sensor data, alerts, recommendations, and logs. Explain the logical links between them and the purpose of these links.

\subsection{Description of Main Database Tables}

This subsection should provide a textual explanation of the main tables and their attributes.

\begin{table}[H]
    \centering
    \caption{Sample table describing important database entities}
    \label{tab:db-entities-placeholder}
    \begin{tabular}{>{\raggedright\arraybackslash}p{0.22\textwidth} >{\raggedright\arraybackslash}p{0.28\textwidth} >{\raggedright\arraybackslash}p{0.36\textwidth}}
        \toprule
        Table & Main fields & Purpose \\
        \midrule
        users & id, email, password\_hash, role & Stores user accounts and access roles \\
        plants & id, user\_id, name, species, location & Stores plant profiles created by users \\
        sensors & id, device\_id, plant\_id, device\_token\_hash & Stores registered sensor devices and bindings \\
        sensor\_data & sensor\_id, plant\_id, moisture, temperature, humidity, light & Stores historical telemetry \\
        alerts & user\_id, plant\_id, severity, status, metric & Stores detected alert conditions \\
        recommendations & plant\_id, alert\_id, text, reason & Stores active care advice \\
        \bottomrule
    \end{tabular}
\end{table}

Table~\ref{tab:db-entities-placeholder} is a sample table block. Replace or extend it with the final table list and descriptions based on the actual schema used in the project.

\section{User Scenarios and Core Functionality}

This section should explain the practical system behavior from the user perspective. Describe how a user interacts with the final prototype and what steps are required to achieve each scenario.

\subsection{User Registration and Plant Profile Management}

Describe how a user creates an account, logs in, and manages plant profiles. Mention fields such as plant name, species, location, and description. Explain why plant profiles are needed to bind telemetry and organize monitoring.

\subsection{Sensor Connection and Plant Monitoring}

Explain how a user creates a sensor record, receives a one-time device token, attaches the sensor to a plant, and then receives readings in the dashboard. Mention both physical device ingest and software simulation as valid demo modes for the diploma presentation.

\subsection{Alerts and AI-Based Recommendations}

Describe how alerts are generated, displayed, and transitioned through their lifecycle. Explain that recommendations are generated as explainable care suggestions and that the dashboard also exposes machine learning support fields such as prediction, confidence, and analysis method.

\section{Implementation and Testing}

This section should describe how the software was implemented and how its correctness was verified.

\subsection{Backend Implementation}

In this subsection, describe the FastAPI backend structure, authentication, business services, repositories, and database interaction. Explain how REST endpoints support user, sensor, monitoring, alert, and admin functionality.

\subsection{Dashboard and Interface Implementation}

Describe the frontend implementation, routes, dashboard widgets, charts, condition cards, alerts page, settings page, and admin view. Mention how the interface displays both traditional analytical outputs and ML-enhanced condition data.

\subsection{Machine Learning Module Implementation}

This subsection should focus on the implemented Random Forest classifier. Explain:
\begin{itemize}
    \item selected input features;
    \item synthetic expert-labeled training dataset;
    \item model training script;
    \item saved model artifact and metadata;
    \item safe runtime loading and fallback behavior;
    \item hybrid integration with rule-based logic.
\end{itemize}

You can cite the theoretical background here again if needed \cite{breiman2001randomforests}.

\subsection{Testing Results}

This subsection should summarize what was tested and what the outcome was. Use repository evidence such as backend tests, frontend build checks, ingest verification, dashboard verification, alert generation, and access-control checks.

Suggested content:
\begin{itemize}
    \item unit tests for security, alert transitions, and condition logic;
    \item machine learning service fallback and prediction tests;
    \item successful frontend build;
    \item successful simulated demo flow;
    \item known warnings or limitations that do not block the defense demo.
\end{itemize}

\section{Limitations and Future Improvements}

This section should honestly describe the current boundaries of the implemented prototype. Mention that the machine learning module is trained on a synthetic expert-labeled dataset and therefore should be presented as a supervised classification support module rather than a clinically or agronomically validated diagnostic engine.

Possible improvement directions:
\begin{itemize}
    \item collecting real sensor datasets for model retraining;
    \item adding species-specific thresholds;
    \item adding richer visualization or mobile support;
    \item improving sensor calibration and reliability;
    \item extending notifications and long-term analytics.
\end{itemize}
